% 核心观点：可靠的语义监督引导任务关系表示，历史条件反馈将新的执行证据转化为槽位对齐的动作修正。
% 叙述逻辑：分割标签与预测质量 -> RGB和语义的联合编码及关系监督 -> 冻结计划上的因果GRU残差。
% 理论对应：伪标签误差与关系读出误差；关系及上下文的保留；新增信息的利用与有效性门控。
% 写作边界：正文描述已有机制；校准、分层审核及可见性状态接入仅作为附录中的扩展协议。
\section{Method}
\label{sec:method}
% 原生矢量机制示意，不是实测架构性能图；坐标不代表实测延迟。
\begin{figure*}[t]
\centering
\begin{tikzpicture}[x=1cm,y=1cm,>=Latex,font=\small,
 block/.style={draw,align=center,minimum height=.65cm,text width=2.8cm},
 flow/.style={->,line width=.5pt}]
\node[block] (rgb) at (1.5,3.2) {多视角RGB与本体};
\node[block] (seg) at (1.5,2.0) {冻结分割器\\软语义图};
\node[block] (enc) at (5.2,3.2) {联合编码器\\RGB + 语义 + 查询};
\node[block] (query) at (9,3.2) {任务关系q-token\\计划时刻$\tau_p$};
\node[block] (act) at (12.8,3.2) {动作解码器\\基础动作块};
\node[block,dashed] (aux) at (9,4.5) {关系辅助头与伪标签\\仅训练监督};
\node[block] (history) at (5.2,.4) {最新RGB、本体\\因果执行历史};
\node[block] (gru) at (9,.4) {有限窗口GRU\\槽位残差预测};
\node[block] (exec) at (12.8,.4) {计划/时效校验\\执行$B_t+\Lambda_t r_t$};
\draw[flow] (rgb)--(enc);
\draw[flow] (rgb.south)--node[right]{RGB}(seg.north);
\draw[flow] (seg.east)-|(enc.south);
\draw[flow] (enc)--(query);
\draw[flow] (enc.north)--(5.2,5.5)--node[above]{完整编码上下文}(12.8,5.5)--(act.north);
\draw[flow,dashed] (query)--(aux);
\draw[flow] (query)--node[right]{缓存查询}(gru);
\draw[flow] (history)--(gru);
\draw[flow] (gru)--(exec);
\draw[flow] (act)--node[right]{基础计划与槽位}(exec);
\draw[flow] (act.south west)--(gru.north east);
\node[align=center,text width=3cm] at (1.5,.4) {分阶段训练：\\分割器 $\to$ 基础策略\\$\to$ 残差分支};
\end{tikzpicture}
\caption{方法机制示意。保留RGB并加入软语义与关系监督；训练残差时冻结分割器及基础策略。
虚线分支是关系辅助监督，不在部署时计算关系控制律。历史记录包含形成时的计划上下文；
反馈使用缓存查询与最新证据预测未来槽位的动作修正。基础和反馈推理共享线程，并非图中两条路径独立并发。}
\label{fig:architecture}
\end{figure*}

\begin{figure}[t]
\centering
\begin{tikzpicture}[x=1cm,y=1cm,>=Latex,font=\scriptsize]
\draw[->] (0,2.0)--(8,2.0) node[below left]{时间};
\node[above] at (1.7,2) {$\tau_p$};
\node[above] at (3.8,2) {$\tau_o$};
\node[above] at (5.8,2) {$t$};
\draw (1.7,1.9)--(1.7,2.1);
\draw (3.8,1.9)--(3.8,2.1);
\draw (5.8,1.9)--(5.8,2.1);
\node[anchor=west] at (0,1.4) {推理线程};
\draw (1.7,1.05) rectangle (3.7,1.7);
\node at (2.7,1.38) {基础更新（优先）};
\draw (3.8,1.05) rectangle (5.3,1.7);
\node at (4.55,1.38) {反馈推理};
\node[anchor=west] at (0,.35) {控制循环};
\draw[->] (1.7,.35)--(7.8,.35);
\node[above] at (3.5,.35) {持续执行缓存计划};
\draw[->] (5.3,1.05)--(5.8,.35);
\node[align=left,anchor=west] at (5.9,.95) {槽位有效\!：采用\\迟到/过期：丢弃};
\end{tikzpicture}
\caption{共享线程的时序示意，非实测延迟。$\tau_p$为缓存计划的观测时刻，$\tau_o$为反馈所用观测时刻，$t$为目标槽位。
基础任务占用期间反馈等待；接收时重新核对计划编号和槽位。同计划重叠候选最后合格者覆盖旧值，不累加。}
\label{fig:scheduling}
\end{figure}

方法围绕三个环节展开：以语义分割和标签质量诊断提供语义信息，
通过语义输入与关系监督引导动作块表示，再利用因果历史对执行中的基础计划进行实时残差修正。
分割与关系监督针对第~\ref{sec:representationanalysis}节的信息保留问题；
历史反馈与槽位对齐针对第~\ref{sec:feedbackanalysis}节的信息利用和执行时效问题。
训练分为分割器、基础策略和残差分支三个阶段，每一阶段冻结已完成的前级模型。
部署时参数固定，控制循环持续执行缓存计划，基础预测与反馈推理异步提交至同一工作线程。
整体数据流见图~\ref{fig:architecture}，共享线程的执行关系见图~\ref{fig:scheduling}。
正文描述已有机制；尚未实现的不确定性评估和标签审核扩展集中列于附录~\ref{app:seguncertainty}，
不作为已有结果的解释依据。

\subsection{语义分割与标签质量处理}
\label{sec:methodseg}
语义监督的作用取决于它是否保留了正确的对象及其关系。
现有流程以视频掩码构造语义监督，并用离线质量诊断降低异常标签对关系学习的影响。

\subsubsection{离线标注与轻量分割}
视频掩码由框、点提示及SAM2传播生成\cite{ravi2024sam2}，
离线脚本以时间IoU、面积和质心变化、连通域异常等定位疑似错误片段。
每个视角使用独立的轻量分割器，以掩码标签训练，采用类别加权交叉熵与前景Dice损失。
类别权重处理类别不平衡，片段质量记录用于重试、标签屏蔽与后续关系监督加权。
策略训练前冻结分割器，并固定类别顺序、图像预处理及模型版本。
离线传播可利用同一示范的后续帧，在线分割只读取当前已获得的图像。
已有残差划分仅保证验证轨迹不参与残差头拟合，冻结基础策略和分割器可能见过这些轨迹。
各阶段实际使用范围须由数据清单核实；全流程独立划分属于后续实验协议，
不追溯性地作为已有模型的训练事实（附录~\ref{app:datascope}）。

\subsubsection{质量加权与有效性}
离线质量分数按相关类别转换为关系监督权重，不直接充当控制门控$\Lambda$。
几何有效性由掩码是否包含相应成员确定，当前面积项的有效标记则恒为真：
空目标掩码产生零面积标签，再乘以质量权重。
因此，已有实现尚不能区分确认不可见、漏标和标注失败；
启发式质量分数不能替代这种判断。
附录~\ref{app:visibilityextension}说明这一实现边界及扩展所需的信息链路。
分层人工抽查、概率校准与关系扰动诊断作为后续评估协议，不进入当前在线路径。

\subsection{任务相关语义的联合编码}
\label{sec:methodrelations}
语义通过两条路径进入动作学习：软语义图提供对象级结构，
关系查询通过辅助监督保留与任务有关的差异。
两条路径均与RGB联合使用，使策略仍能利用语义类别未覆盖的姿态、纹理和接触线索。
同一图像经冻结分割器得到的语义提供归纳引导，而不增加物理观测；
其目的在于帮助联合表示保留与控制后果有关的信息\cite{zhang2021dbc,subramanian2022ais}。

\subsubsection{RGB与软语义输入}
视角$v$的分割器根据图像$I_t^v$输出类别概率$p_t^v(c,\xi)$，
并通过固定颜色向量$e_c^{\mathrm{sem}}$形成三通道软语义图：
\begin{equation}
 S_t^v(\xi)=\sum_{c\in\mathcal C_v}p_t^v(c,\xi)e_c^{\mathrm{sem}}.
 \label{eq:methodsoftsem}
\end{equation}
其中，$\mathcal C_v$为该视角的类别集合，$\xi$为像素位置，
$e_c^{\mathrm{sem}}\in[0,1]^3$为类别$c$的固定颜色编码。
各视角维护明确的类别对应关系；软颜色图是概率的压缩表示，不采用argmax硬化，
也不等同于完整的多类别概率张量。

每个视角的RGB与软语义图分别经共享视觉骨干和投影变成视觉token，
与本体信息共同输入ACT编码器\cite{zhao2023act}。
多个视角作为独立视觉流联合编码，不在像素坐标上直接平均掩码或质心。
对尚未标定到共同坐标系的关系，保留视角索引；
跨视角信息的结合由联合编码器学习。

\subsubsection{关系监督查询}
根据给定任务预先选择$K_{\mathrm{rel}}$组关系$G^{(j)}\in\R^{d_j}$，
例如可见程度、相对位置或区域占据关系，并规定其坐标、尺度与有效条件。
各组共同组成第~\ref{sec:representationanalysis}节的关系向量$G$，
其中$d_j$为第$j$组的维数，$\sum_jd_j=d_G$。
离线掩码$\widetilde M_t^{1:V}$经固定算子$\rho_j$生成关系伪标签
$\widetilde G_t^{(j)}=\rho_j(\widetilde M_t^{1:V})$。
这些算子只用于构造监督，不向部署策略提供真实物体状态。
具体对象、关系数量和计算方式作为任务实例给出，而不固定在通用架构中。

每组关系对应一个可学习查询token，与视觉和本体token共同进入编码器。
其输出$z_t^{(j)}$由辅助头读出关系预测
$\widehat G_t^{(j)}=D_j(z_t^{(j)})$，
同时作为编码器上下文供动作解码器读取。
因此，关系通过受监督的隐藏表示影响动作块，
而不是由人工规则将某个关系数值直接转换为动作或阶段切换。

训练使用有效性与质量加权的关系损失，并与ACT动作目标联合优化：
\begin{align}
 \mathcal L_{\mathrm{rel}}
 &=\frac{1}{K_{\mathrm{rel}}}\sum_{j=1}^{K_{\mathrm{rel}}}
 \frac{\sum_{n,c}w_{njc}\,
 \ell_{\beta_G}(\widehat G_{njc}-\widetilde G_{njc})}
 {\max(1,\sum_{n,c}w_{njc})},\nonumber\\
 \mathcal L_{\mathrm{base}}
 &=\mathcal L_{\mathrm{act}}
   +\lambda_{\mathrm{KL}}\mathcal L_{\mathrm{KL}}
   +\lambda_{\mathrm{rel}}\mathcal L_{\mathrm{rel}}.
 \label{eq:methodbaseloss}
\end{align}
其中，$n$为批内样本索引，$c$为组内分量索引，$w_{njc}\in[0,1]$为有效性与质量权重，
$\ell_{\beta_G}$为转折参数$\beta_G>0$的Smooth L1损失；
$\mathcal L_{\mathrm{act}}$为有效动作槽位上的L1模仿损失，
$\mathcal L_{\mathrm{KL}}$为ACT变分潜变量的KL正则，
$\lambda_{\mathrm{KL}},\lambda_{\mathrm{rel}}\ge0$为损失权重。
各关系组先按固定尺度归一化，未定义关系的权重置零。
关系标签不可用时保留该样本的动作监督，避免丢失遮挡等困难情形。
分割器在此阶段保持冻结。

ACT训练使用由本体观测和示范动作块编码、采样得到的CVAE潜变量；
该潜变量与关系查询共同进入Transformer编码器，因此关系辅助损失也在这一条件下计算。
部署沿用ACT将潜变量置零的约定\cite{zhao2023act}；
残差缓存同样以基础策略的部署模式生成，不输入示范动作。
后续关系读出评价须使用置零潜变量的模式，训练关系损失不能代替该评价。

这一设计分别作用于式\eqref{eq:pseudolabel}中的监督来源与关系读出误差。
所选关系是否足以解释决策仍取决于任务及上下文；
因此，后续反馈保留原始视觉、本体与执行历史，而非仅以这些关系数值作为策略状态。

\subsection{历史条件实时残差反馈}
\label{sec:methodfeedback}
基础动作块提供连续的行为计划，反馈分支用新观测与此前执行响应修正未来槽位。
控制循环持续执行缓存计划；基础预测与反馈推理共享一个工作线程和一个待完成任务，
空闲时优先提交到期的基础更新，否则提交到期的反馈任务。
因此，两类推理并不独立并发，基础推理期间反馈也会等待；
实际更新间隔由资源占用与结果有效性共同决定。
残差始终相对于当前有效计划定义，基础策略在残差训练和部署中均保持冻结。

\subsubsection{计划上下文与因果GRU}
基础策略在观测时刻$\tau_p$生成动作块及关系查询上下文。
反馈在时刻$\tau_o$获取新图像，为尚未执行的槽位$t$生成候选修正，
沿用$\tau_p\le\tau_o\le t$的时刻约定。
查询在两次计划更新之间缓存，新图像则在每次反馈更新时重新编码。

第$k$次反馈更新的输入记录$e_k^{\mathrm{fb}}$联合编码新图像特征、
缓存查询、本体状态、因果速度估计、上一发送指令，
以及当前计划中的未来基础动作、槽位有效性和时序信息。
各历史记录保留其形成时的计划上下文，不以后来生成的计划替换。
对最近$L_{\mathrm{hist}}$条记录，用固定窗口GRU获得反馈隐藏状态：
\begin{equation}
 h_i=\operatorname{GRU}_{\omega}(e_i^{\mathrm{fb}},h_{i-1}),
 \quad i=\max(1,k-L_{\mathrm{hist}}+1),\ldots,k.
 \label{eq:methodgru}
\end{equation}
其中，$\omega$为GRU参数，$L_{\mathrm{hist}}$为历史窗口长度；
每次计算将窗口起点之前的隐藏状态置零，并取最后一条有效记录的输出$h_k$。
旧图像特征使用缓存，断流或重置时清空历史。
由此，GRU可以结合此前发出的指令与随后观察到的变化解释当前画面，
而不只对相邻动作作平滑。

缓存查询描述$\tau_p$的任务关系，最新图像与历史补充计划生成后的变化。
这些输入及运行记录共同构成理论中的$Z_t^{\mathrm{fb}}$。
其对执行时刻关系$G_t$的保留能力需单独检验，
不能以基础查询在$\tau_p$的辅助预测精度代替。
实验将采用关系读出探针与相同当前信息的MLP对照检验这一联系，
不额外假定GRU已经恢复了完整状态。

\subsubsection{残差预测与监督}
反馈头根据$h_k$输出$H_r$个未来槽位的残差，
第$j$个槽位的候选指令由基础动作$B_{t_{k,j}}$与限幅残差组成：
\begin{equation}
 \begin{aligned}
 r_{k,j}&=\bm\beta\odot\tanh([W_rh_k+b_r]_j),\\
 \widehat u_{k,j}&=B_{t_{k,j}}+r_{k,j},
 \qquad j=0,\ldots,H_r-1.
 \end{aligned}
 \label{eq:methodresidual}
\end{equation}
其中，$t_{k,j}$为预定目标执行槽位，$\bm\beta\in\R_{>0}^{d_u}$为逐动作维度限幅，
$W_r,b_r$为输出头参数，$\odot$为逐元素乘积。
输出头零初始化，限幅仅由训练数据和平台动作单位确定。
这里的$r_{k,j}$是动作指令残差，不是语义状态增量，也不在相邻时刻累积积分。

残差训练沿示范轨迹回放冻结基础策略，以固定采样步长和预设推理延迟生成计划与目标槽位。
缓存遵循同一槽位定义，但不模拟在线共享工作线程的资源竞争及不规则更新间隔。
未来示范动作仅作为监督目标，不进入反馈输入。
对有效槽位，使用尺度归一化的纠偏动作误差及残差正则：
\begin{equation}
 \begin{aligned}
 \mathcal L_{\mathrm{res}}
 &=\frac{1}{d_u\max(1,\sum_{n,j}m_{nj})}
 \sum_{n,j,c}m_{nj}\\
 &\quad{}\times\Big[\ell_{\beta_r}\!\left(
 \frac{B_{njc}+r_{njc}-u^{\mathrm E}_{njc}}{\sigma_c}
 \right)\\
 &\qquad{}+\lambda_r(r_{njc}/\sigma_c)^2\Big].
 \end{aligned}
 \label{eq:methodresloss}
\end{equation}
其中，$n$为回放样本索引，$j$为目标槽位索引，$c$为动作维度索引，
$m_{nj}$为槽位有效标记，$u^{\mathrm E}_{njc}$为对应示范指令，
$\sigma_c>0$为训练集动作尺度，$\beta_r>0$与$\lambda_r\ge0$为损失参数。
回放输入来自示范记录的状态与因果指令历史，未模拟残差造成的新环境响应；
因此，离线拟合之外仍需用自主闭环评估检验部署效果。

\subsubsection{目标槽位对齐与回退}
反馈输出绑定计划编号及绝对目标槽位。
预定的推理延迟决定首个目标槽位，输出区间按目标反馈间隔配置；
资源等待可能使实际修正覆盖出现间隙。时间戳到控制步的映射见附录~\ref{app:slottiming}。
同一计划的重叠槽位由最后通过接收校验的候选覆盖，
新计划生效时清空旧计划修正；迟到结果只丢弃过期前缀，不将其平移到后续时间。
这些规则确保纠偏针对生成时所条件化的基础动作，而不叠加到已经修正的指令上。

执行时检查计划归属、槽位、数值有效性及观测和候选年龄，
按第~\ref{sec:feedbackanalysis}节的规则执行$u_t=B_t+\Lambda_t r_t$。
其中，$r_t$为确定性候选选择后的残差，$\Lambda_t$为执行有效性门控。
无有效修正时执行有效基础动作；启动等待首个计划时保持初始化指令；
基础计划失效时终止当前策略执行并进入平台预先规定的处理流程。
该回退同时纳入参考策略$\pi_0$，关闭残差，而不无限沿用过期动作。

总体上，关系监督促使基础表示关注任务相关差异，
GRU尝试利用动作--响应历史，槽位校验保证候选作用于预定执行位置。
这些训练与执行机制不直接优化理论中的$Q_t^0$；
依据式\eqref{eq:closedlooptaskgain}，最终贡献以闭环任务结果检验，
并用关系误差、当前信息对照及逐槽延迟日志解释。
具体任务定义、实现设置及待完成部分列于附录~\ref{app:methoddetails}。
